% !TeX spellcheck = ru_RU
% !TEX root = vkr.tex

\section*{Введение}
\thispagestyle{withCompileDate}

Конечные смеси вероятностных распределений используются для описания наблюдений, которые формируются несколькими скрытыми источниками. Каждый источник представляется компонентой с собственным распределением, а его доля --- весом компоненты. В этой работе рассматриваются одномерные непрерывные наблюдения. Компоненты могут принадлежать одному семейству либо разным семействам, например нормальному и экспоненциальному~\cite{McLachlanPeel2000}. Такие модели позволяют описывать асимметричные плотности, перекрывающиеся группы и небольшие доли наблюдений с отличающимся поведением.

При заданных числе компонент и их семействах требуется оценить параметры и веса по выборке. Для этого применяются метод ожидания и максимизации, его варианты с условным оцениванием групп параметров, прямая численная оптимизация и методы моментов. Аналитические формулы доступны только для определённых семейств и ограничений; в остальных случаях требуется численное решение. Результат также зависит от начального приближения, а число компонент и их семейства могут быть неизвестны. Поэтому программные средства должны разделять оценивание при заданной структуре и выбор самой структуры.

Работа выполняется в рамках развития \texttt{pysatl-mpest} --- библиотеки Python для оценивания параметров смесей. В проекте уже имеются общее представление смеси, интерфейсы оценщиков и отдельные итерационные стратегии. Вычислительное ядро \texttt{pysatl-core} предоставляет средства работы с распределениями и параметрическими семействами~\cite{MpestRepository,CoreRepository}. Развитие библиотеки требует согласовать оценщики с ядром, уточнить управление свободными и фиксированными параметрами, выделить построение начальных приближений и определить проверяемую процедуру выбора структуры.

Для выбора структуры используются критерии, которые сопоставляют качество подгонки и число свободных параметров. Дополнительную проверку добавления компоненты можно строить параметрическим бутстрепом: из оценённой более простой модели генерируются повторные выборки, на которых сравниваются обе модели~\cite{McLachlan1987}. Для проверки самой библиотеки выбираются положительные смеси Вейбулла, разнородная логнормально-паретовская смесь и реальные страховые убытки. На синтетических данных известны истинные параметры; на реальных данных проверяется качество подгонки.

Близкие программные решения уже поддерживают оценивание смесей и расширение компонентных моделей. Поэтому актуальность работы связана с конкретным развитием библиотеки проекта и воспроизводимой проверкой её стратегий. Важен вопрос, при каких условиях общий механизм численной оптимизации и специализированные аналитические обновления дают согласованные результаты и как различаются их вычислительные затраты. Второй вопрос касается выбора числа компонент и семейств при ограниченном бюджете оценивания. Ответы должны опираться на сравнение методов и эксперименты с явно заданными условиями.

На текущем этапе подготовлен обзор математических методов, программных инструментов и прикладных сценариев. На его основе определяются архитектурные границы, приоритетная функциональность и фрагменты опубликованных исследований для будущего воспроизведения. Планируемые результаты разработки и экспериментов отделены от результатов обзора.
